% =========================================================
% Compléments M1 — Variation quadratique (C. Dérivations)
% =========================================================

\subsection*{(1) De la convergence en $L^2$ à la convergence p.s. : une esquisse propre}
Dans le texte principal, on a montré sur une partition uniforme que
\[
Q_n=\sum_{k=0}^{n-1}(\Delta W_k)^2 \to T \quad \text{en }L^2.
\]
Pour obtenir une convergence \emph{presque sûre}, on peut utiliser une sous-suite dyadique.
Posons $n_m=2^m$ et $\Delta t_m=T/2^m$, et $Q_{n_m}$ la somme correspondante.
On a (calcul déjà fait) :
\[
\Var(Q_{n_m})=2T\Delta t_m = 2T\frac{T}{2^m}=\frac{2T^2}{2^m}.
\]
Par l'inégalité de Tchebychev, pour tout $\varepsilon>0$,
\[
\Prob(|Q_{n_m}-T|>\varepsilon)\le \frac{\Var(Q_{n_m})}{\varepsilon^2}
=\frac{2T^2}{\varepsilon^2}\,2^{-m}.
\]
La série $\sum_m \Prob(|Q_{n_m}-T|>\varepsilon)$ converge, donc par Borel--Cantelli :
\[
Q_{n_m}\to T \quad \text{p.s.}
\]
\paragraph{Remarque (niveau M1).}
Passer ensuite de la sous-suite dyadique à une suite générale demande un argument de stabilité (contrôle des raffinements). L'idée est que raffiner une partition n'ajoute qu'un terme dont l'espérance est contrôlée et dont la variance décroît avec le pas maximal.

\subsection*{(2) Covariation : calcul explicite en dimension 2}
Soient $W^{(1)}$ et $W^{(2)}$ deux browniens corrélés tels que
\[
\dd W_t^{(2)}=\rho\,\dd W_t^{(1)}+\sqrt{1-\rho^2}\,\dd B_t,
\]
où $B$ est un brownien indépendant de $W^{(1)}$.
Alors, en utilisant la table d'Itô :
\[
\dd[W^{(1)},W^{(2)}]_t=\dd W_t^{(1)}\,\dd W_t^{(2)}
=\rho(\dd W_t^{(1)})^2+\sqrt{1-\rho^2}\,\dd W_t^{(1)}\dd B_t
=\rho\,\dd t.
\]
En intégrant :
\[
\boxed{[W^{(1)},W^{(2)}]_t=\rho t.}
\]
Ce calcul est exactement ce qui se cache derrière la phrase : «corrélation instantanée $\rho$».

\subsection*{(3) Crochet d'un processus d'Itô : dérivation pas à pas}
Soit
\[
X_t=X_0+\int_0^t a_s\,\dd s+\int_0^t b_s\,\dd W_s.
\]
On note $A_t=\int_0^t a_s\,\dd s$ (variation finie) et $M_t=\int_0^t b_s\,\dd W_s$ (martingale continue).
Alors $X=A+M$ et
\[
[X]_t=[A+M]_t=[A]_t+[M]_t+2[A,M]_t.
\]
Or $[A]_t=0$ (variation finie), et $[A,M]_t=0$ (produit d'une variation finie et d'une martingale continue), donc
\[
[X]_t=[M]_t.
\]
Il reste à calculer $[M]_t$.
Sur une partition, $\Delta M_k\approx b_{t_k}\Delta W_k$, donc heuristiquement
$\sum (\Delta M_k)^2\approx \sum b_{t_k}^2(\Delta W_k)^2\to \int_0^t b_s^2\,\dd s$.
La preuve rigoureuse passe par approximation de $b$ par des processus simples et utilisation de l'isométrie d'Itô.
On obtient :
\[
\boxed{[X]_t=\int_0^t b_s^2\,\dd s.}
\]

\subsection*{(4) Exemple : variation quadratique d'un GBM et variance réalisée}
Sous GBM, $\dd S_t=(r-q)S_t\,\dd t+\sigma S_t\,\dd W_t$.
La partie martingale est $\int_0^t \sigma S_s\,\dd W_s$, donc
\[
[S]_t=\int_0^t \sigma^2 S_s^2\,\dd s.
\]
\paragraph{Lien finance.}
Si l'on observe des prix à haute fréquence et qu'on somme les carrés des rendements (réalisée),
on approxime une quantité de la forme $\int_0^t \sigma^2 S_s^2\,\dd s$ (ou, en log, $\int_0^t \sigma^2\,\dd s$).
C'est l'origine théorique des estimateurs de volatilité réalisée.

\subsection*{(5) Exercices corrigés (variation quadratique)}
\paragraph{Exercice 1 --- $W_t^2-t$ est une martingale.}
Montrer que $M_t=W_t^2-t$ est une martingale et en déduire que $[W]_t=t$.

\emph{Solution.}
Appliquer Itô à $f(x)=x^2$ :
\[
\dd(W_t^2)=2W_t\,\dd W_t+(\dd W_t)^2=2W_t\,\dd W_t+\dd t.
\]
Donc $\dd(W_t^2-t)=2W_t\,\dd W_t$ : c'est une intégrale d'Itô, donc une martingale.
On lit aussi que le terme correctif est précisément $(\dd W_t)^2=\dd t$, ce qui correspond à $[W]_t=t$.

\paragraph{Exercice 2 --- covariation linéaire.}
Soit $X_t=a t+b W_t$ avec $a,b\in\R$. Calculer $[X]_t$.

\emph{Solution.}
La partie variation finie $at$ a crochet nul, donc
\[
[X]_t=[bW]_t=b^2[W]_t=b^2 t.
\]

\paragraph{Exercice 3 --- crochet d'une intégrale.}
Soit $M_t=\int_0^t H_s\,\dd W_s$ avec $H\in\mathcal{H}^2$. Montrer que $[M]_t=\int_0^t H_s^2\,\dd s$.

\emph{Solution (esquisse).}
Pour $H$ simple, $M$ est une somme $\sum H_{t_k}\Delta W_k$, donc la somme des carrés $\sum (\Delta M_k)^2$ est une somme de $H_{t_k}^2(\Delta W_k)^2$ qui converge vers $\int H^2$.
On étend par densité aux $H$ généraux via l'isométrie d'Itô.
